%% AffectAI Dataset Paper — NeurIPS 2026 Evaluations & Datasets Track
%%
%% Track: https://neurips.cc/Conferences/2026/CallForEvaluationsDatasets
%% Deadlines: abstract May 4 2026 AoE; full paper May 6 2026 AoE
%% Format: NeurIPS 2026 LaTeX template
%%
%% Compile: pdflatex main → bibtex main → pdflatex × 2
%%
%% Review mode: DOUBLE-BLIND.
%% For submission, use the official NeurIPS 2026 style in anonymous mode.
%% At camera-ready: switch to [final], restore authors/affiliations/ORCIDs,
%% acknowledgements, and de-anonymise dataset URLs.

\documentclass{article}

%% NOTE:
%% For a double-blind NeurIPS submission, the usual setting is:
%%   \usepackage{neurips_2026}
%% not [preprint].
%% Use [preprint] only for arXiv or internal public drafts.
\usepackage{neurips_2026}
% \usepackage[preprint]{neurips_2026}  % use only for non-anonymous preprint
% \usepackage[final]{neurips_2026}     % use at camera-ready

\usepackage[T1]{fontenc}
\usepackage[utf8]{inputenc}
\usepackage{microtype}
\usepackage{booktabs}
\usepackage{multirow}
\usepackage{graphicx}
\usepackage{subcaption}
\usepackage{xcolor}
\usepackage{url}
\usepackage{hyperref}
\usepackage{cleveref}

%% Title — keep concise and descriptive.
\title{AffectAI: A Multimodal Dataset of Co-Located Group Interaction}

%% Anonymous author block for double-blind review.
\author{%
  Anonymous Authors\\
  Submission to NeurIPS 2026 Evaluations \& Datasets Track\\
  \texttt{\{anon@example.com\}}%
}

\begin{document}
\maketitle

%% ═══════════════════════════════════════════════════════════════════════════
\begin{abstract}
We present \textbf{AffectAI}, a multimodal dataset of co-located
four-person group interaction during structured collaborative tasks:
information pooling, mini-negotiation, idea generation, and a public-goods
micro-game.
Each session captures per-participant physiological signals, 2D eye tracking,
close-talk audio, tablet-based self-reports, and pre-session Big-Five
personality scores, linked through anonymised seat identifiers.
The paper focuses on dataset characterisation rather than leaderboard-style
benchmarking: we document the study design, synchronisation strategy,
modality coverage, missingness, task outcomes, and release structure.
Where feature availability permits, we include one conservative worked example
illustrating how task-level multimodal summaries can be analysed; otherwise,
the section specifies supported analyses without making strong predictive
claims.
The release is organised in a BIDS-inspired structure and accompanied by
Croissant metadata, datasheet documentation, quality-control reports, and
processing scripts.
We explicitly document the limitations of the current release, including
small sample size, single-site collection, single-language interaction,
per-modality missingness, 2D rather than world-frame gaze, and voice
re-identification risks.
Dominance-related analyses, video-based 3D pose, and world-frame gaze are
reserved for follow-up work.
\end{abstract}

%% ═══════════════════════════════════════════════════════════════════════════
\section{Introduction}
\label{sec:intro}

Most multimodal corpora used to study affect and interaction have been
recorded in single-participant or dyadic settings
\cite{ringeval2013introducing,busso2008iemocap,koelstra2012deap}.
Co-located groups impose different constraints: turn-taking is multi-party,
gaze targets multiply, roles can shift within a session, and physiological
responses must be interpreted against an evolving social context rather than
a fixed stimulus.
Datasets that capture this setting are relatively rare; datasets that combine
per-participant physiology, gaze, audio, task outcomes, self-reports, and
personality measures in the same co-located group setting are rarer still.

We present \textbf{AffectAI}, a multimodal corpus of four-person group
interaction designed to support analysis of affective, cognitive, and
conversational dynamics in structured collaborative tasks.
Each session includes four task blocks chosen to elicit complementary
social-affective situations: hidden-profile information pooling,
mini-negotiation, nominal-group-technique idea generation, and a public-goods
micro-game.
For each participant, we record synchronised physiology, 2D eye tracking,
close-talk audio, and tablet-based self-reports.
We additionally administer the Big-Five Inventory before each session.

The current paper is deliberately framed as a dataset-characterisation paper.
Its primary contribution is not a new model or a leaderboard, but a transparent
description of what the dataset contains, how the modalities are synchronised,
what can be evaluated with the current release, and where the main limitations
lie.
This framing is appropriate for the current scale of the dataset: the corpus is
rich and multimodal, but the number of groups is too small to support broad
benchmark claims.

The paper includes one worked example if the required feature tables are ready
and reliable by the internal freeze.
This worked example is intended to demonstrate dataset use, not to establish a
definitive scientific result.
If one or more feature families are incomplete, the worked example is reduced
to the available subset, and the missing modality is reported as a limitation.

The contributions of this paper are:
\begin{enumerate}
  \item A multimodal corpus of \textbf{TODO~$N_g$} groups
    (\textbf{TODO~$N_s$} analysable sessions;
    \textbf{TODO~$N_p$} unique participants;
    \textbf{TODO~$T$} hours total) capturing per-participant physiology,
    2D gaze, close-talk audio, tablet self-reports, task outcomes, and
    Big-Five personality scores.
  \item A dataset-characterisation analysis covering synchronisation,
    modality coverage, per-session missingness, task structure, and
    quality-control outputs.
  \item One conservative worked example, contingent on final feature
    availability, together with formal specifications for additional analyses
    that the dataset supports.
  \item A reproducible processing and release structure, including
    BIDS-inspired organisation, Croissant metadata, datasheet documentation,
    and transparent access-tier decisions for sensitive modalities.
\end{enumerate}

The remainder of the paper is organised as follows.
\Cref{sec:related} surveys related corpora.
\Cref{sec:design} describes the task and session design.
\Cref{sec:modalities} details the acquisition setup.
\Cref{sec:processing} describes processing and release structure.
\Cref{sec:worked} presents the worked example and supported analyses.
\Cref{sec:stats} reports coverage and quality.
\Cref{sec:limitations} states limitations.
\Cref{sec:ethics} discusses ethics, access, and responsible use.

%% ═══════════════════════════════════════════════════════════════════════════
\section{Related Work}
\label{sec:related}

% TODO ~3/4 page.
% Position against:
%   - Group multimodal corpora: AMI, ICSI, ELEA, MatchNMingle, UDIVA.
%   - Affect + physiology corpora: RECOLA, K-EmoCon, MAHNOB-HCI, DEAP.
%   - Conversational corpora with transcripts/diarisation: AMI, ICSI, CHiME.
%
% Suggested comparison axes:
%   group_size / co_located / structured_tasks / per_person_physio /
%   per_person_gaze / per_person_audio / self_report / personality /
%   public_or_gated_release.
%
% Core positioning:
% AffectAI is not larger than prior corpora. Its contribution is the
% simultaneous combination of modalities in a structured co-located group
% setting with per-person sensing and transparent release documentation.

\textbf{TODO} short prose overview, followed by a comparison table.
% \input{tables/related_axes}

%% ═══════════════════════════════════════════════════════════════════════════
\section{Dataset Design}
\label{sec:design}

\paragraph{Participants.}
Data were collected from \textbf{TODO~$N_g$} groups of four participants each
($N_p=\textbf{TODO}$ unique participants;
mean age~=~\textbf{TODO};
sex/gender distribution~=~\textbf{TODO}).
Recruitment used a convenience sample drawn from
\textbf{[ANON --- recruitment context]}.
The demographic consequences of this sampling strategy are discussed in
\cref{sec:limitations}.
Each participant gave written informed consent under approval
\textbf{[ANON --- ethics reference]}.
Within each session, participants are referred to using anonymous seat
identifiers P1--P4.
The mapping from real identity to seat identifier is not part of the release.

\paragraph{Pre-session questionnaire.}
Before each session, participants completed the Big-Five Inventory
(BFI-44, \cite{john1999bigfive}).
Per-domain scores are computed as item means after reverse coding, following
standard scoring.
Of \textbf{TODO} unique participants, \textbf{TODO} have matched BFI-44
responses.
The matching and audit procedure are documented in the appendix.

\paragraph{Session structure.}
Each session lasted approximately \textbf{TODO} minutes and comprised a
baseline phase followed by four collaborative tasks.
Task phases are demarcated by experimenter-issued LSL markers and written
into release-side task-window files so downstream users do not need to
reconstruct the timing manually.

\paragraph{Tasks.}
The four tasks were selected to elicit distinct interaction patterns while
remaining understandable to participants and feasible within a single session.

\textbf{T1} is a hidden-profile decision task in which information is
unevenly distributed across participants; successful performance requires
information pooling.

\textbf{T2} is a mini-negotiation task in which participants or subgroups
hold partially conflicting objectives.

\textbf{T3} is a nominal-group-technique idea-generation task combining
individual ideation and group ranking.

\textbf{T4} is a public-goods micro-game with private contribution decisions
and a public reveal phase.

The per-task outcomes are recorded in
\texttt{beh/<session>\_stimuli\_answers.tsv} where available.

\paragraph{Self-report probes.}
At scheduled moments during the session, participants respond to
tablet-based self-report probes.
The collected scales include valence, arousal, and dominance-style ratings.
For this paper, dominance-related analysis is deliberately not treated as a
central contribution; dominance is retained as a released variable where
available and reserved for follow-up work.
The current analysis focuses on dataset characterisation and, where possible,
valence/arousal or task-level summaries.

%% Table: task_overview.tex
%% %% Table: task_overview.tex
%% %% Table: task_overview.tex
%% \input{tables/task_overview} from main.tex
%% Ported from ICMI version with minor cleanup. NeurIPS template uses booktabs.

\begin{table}[t]
  \centering
  \caption{Session structure: task phases with approximate durations and
    primary social/affective dynamics elicited.}
  \label{tab:tasks}
  \small
  \begin{tabular}{llp{4cm}r}
    \toprule
    \textbf{ID} & \textbf{Task} & \textbf{Social dynamic} & \textbf{Duration (s)} \\
    \midrule
    T0 & Baseline / introduction       & Familiarisation, free conversation              & $\sim$300 \\
    T1 & Hidden-Profile decision       & Information asymmetry, persuasion, pooling      & $\sim$480 \\
    T2 & Mini-negotiation              & Conflict resolution, compromise                 & $\sim$540 \\
    T3 & Idea generation (NGT)         & Solitary ideation alternating with group ranking & $\sim$630 \\
    T4 & Public-goods micro-game       & Cooperation, trust, fairness reveal             & $\sim$360 \\
    \bottomrule
  \end{tabular}
\end{table}
 from main.tex
%% Ported from ICMI version with minor cleanup. NeurIPS template uses booktabs.

\begin{table}[t]
  \centering
  \caption{Session structure: task phases with approximate durations and
    primary social/affective dynamics elicited.}
  \label{tab:tasks}
  \small
  \begin{tabular}{llp{4cm}r}
    \toprule
    \textbf{ID} & \textbf{Task} & \textbf{Social dynamic} & \textbf{Duration (s)} \\
    \midrule
    T0 & Baseline / introduction       & Familiarisation, free conversation              & $\sim$300 \\
    T1 & Hidden-Profile decision       & Information asymmetry, persuasion, pooling      & $\sim$480 \\
    T2 & Mini-negotiation              & Conflict resolution, compromise                 & $\sim$540 \\
    T3 & Idea generation (NGT)         & Solitary ideation alternating with group ranking & $\sim$630 \\
    T4 & Public-goods micro-game       & Cooperation, trust, fairness reveal             & $\sim$360 \\
    \bottomrule
  \end{tabular}
\end{table}
 from main.tex
%% Ported from ICMI version with minor cleanup. NeurIPS template uses booktabs.

\begin{table}[t]
  \centering
  \caption{Session structure: task phases with approximate durations and
    primary social/affective dynamics elicited.}
  \label{tab:tasks}
  \small
  \begin{tabular}{llp{4cm}r}
    \toprule
    \textbf{ID} & \textbf{Task} & \textbf{Social dynamic} & \textbf{Duration (s)} \\
    \midrule
    T0 & Baseline / introduction       & Familiarisation, free conversation              & $\sim$300 \\
    T1 & Hidden-Profile decision       & Information asymmetry, persuasion, pooling      & $\sim$480 \\
    T2 & Mini-negotiation              & Conflict resolution, compromise                 & $\sim$540 \\
    T3 & Idea generation (NGT)         & Solitary ideation alternating with group ranking & $\sim$630 \\
    T4 & Public-goods micro-game       & Cooperation, trust, fairness reveal             & $\sim$360 \\
    \bottomrule
  \end{tabular}
\end{table}


%% ═══════════════════════════════════════════════════════════════════════════
\section{Modalities and Acquisition}
\label{sec:modalities}

This release focuses on per-participant streams and task/self-report
annotations.
Multi-camera video streams were recorded for separate follow-up work, but
video-based 3D pose, world-frame gaze, and related analyses are not included
in the present submission.

\Cref{tab:modalities} summarises the release modalities.

%% Table: modalities.tex
%% %% Table: modalities.tex
%% %% Table: modalities.tex
%% \input{tables/modalities} from main.tex
%%
%% This version reflects the modalities released in this paper. Multi-camera
%% video streams are deferred to follow-up work and are not analysed here.

\begin{table}[t]
  \centering
  \caption{Modalities released in this version of AffectAI.
    All per-participant streams use seat identifiers P1--P4.
    Sampling rates marked with $\dagger$ are channel-dependent;
    see the per-channel table in the appendix.
    Multi-camera video streams were recorded for follow-up work; they are not part of the current release or analysis.}
  \label{tab:modalities}
  \small
  \begin{tabular}{llllr}
    \toprule
    \textbf{Modality} & \textbf{Device} & \textbf{Signals} & \textbf{Rate} & \textbf{Per session} \\
    \midrule
    Physiology
      & EmotiBit
      & PPG (3 wavelengths), EDA, temperature, IMU
      & $\dagger$
      & 4 (one per P) \\
    \addlinespace
    Eye tracking (2D)
      & Tobii Pro Glasses 3
      & 2D gaze, pupil diameter, validity
      & 50~Hz
      & 4 (one per P) \\
    \addlinespace
    Audio (close-talk)
      & DPA 4060
      & Mono, 48~kHz / 24-bit
      & 48~kHz
      & 4 (one per P) \\
    \addlinespace
    Audio (room)
      & DPA 4060
      & Mono, 48~kHz / 24-bit
      & 48~kHz
      & 1 \\
    \addlinespace
    Behavioural self-report
      & Android tablets
      & VAD probes, task responses
      & event-driven
      & 4 (one per P) \\
    \addlinespace
    Personality (pre-session)
      & Online questionnaire
      & BFI-44 item responses, demographics
      & one-shot
      & 4 (one per P) \\
    \addlinespace
    Conversation transcripts
      & WhisperX (per-channel)
      & Word-level diarised transcripts
      & event-driven
      & 1 (merged) \\
    \bottomrule
  \end{tabular}
\end{table}
 from main.tex
%%
%% This version reflects the modalities released in this paper. Multi-camera
%% video streams are deferred to follow-up work and are not analysed here.

\begin{table}[t]
  \centering
  \caption{Modalities released in this version of AffectAI.
    All per-participant streams use seat identifiers P1--P4.
    Sampling rates marked with $\dagger$ are channel-dependent;
    see the per-channel table in the appendix.
    Multi-camera video streams were recorded for follow-up work; they are not part of the current release or analysis.}
  \label{tab:modalities}
  \small
  \begin{tabular}{llllr}
    \toprule
    \textbf{Modality} & \textbf{Device} & \textbf{Signals} & \textbf{Rate} & \textbf{Per session} \\
    \midrule
    Physiology
      & EmotiBit
      & PPG (3 wavelengths), EDA, temperature, IMU
      & $\dagger$
      & 4 (one per P) \\
    \addlinespace
    Eye tracking (2D)
      & Tobii Pro Glasses 3
      & 2D gaze, pupil diameter, validity
      & 50~Hz
      & 4 (one per P) \\
    \addlinespace
    Audio (close-talk)
      & DPA 4060
      & Mono, 48~kHz / 24-bit
      & 48~kHz
      & 4 (one per P) \\
    \addlinespace
    Audio (room)
      & DPA 4060
      & Mono, 48~kHz / 24-bit
      & 48~kHz
      & 1 \\
    \addlinespace
    Behavioural self-report
      & Android tablets
      & VAD probes, task responses
      & event-driven
      & 4 (one per P) \\
    \addlinespace
    Personality (pre-session)
      & Online questionnaire
      & BFI-44 item responses, demographics
      & one-shot
      & 4 (one per P) \\
    \addlinespace
    Conversation transcripts
      & WhisperX (per-channel)
      & Word-level diarised transcripts
      & event-driven
      & 1 (merged) \\
    \bottomrule
  \end{tabular}
\end{table}
 from main.tex
%%
%% This version reflects the modalities released in this paper. Multi-camera
%% video streams are deferred to follow-up work and are not analysed here.

\begin{table}[t]
  \centering
  \caption{Modalities released in this version of AffectAI.
    All per-participant streams use seat identifiers P1--P4.
    Sampling rates marked with $\dagger$ are channel-dependent;
    see the per-channel table in the appendix.
    Multi-camera video streams were recorded for follow-up work; they are not part of the current release or analysis.}
  \label{tab:modalities}
  \small
  \begin{tabular}{llllr}
    \toprule
    \textbf{Modality} & \textbf{Device} & \textbf{Signals} & \textbf{Rate} & \textbf{Per session} \\
    \midrule
    Physiology
      & EmotiBit
      & PPG (3 wavelengths), EDA, temperature, IMU
      & $\dagger$
      & 4 (one per P) \\
    \addlinespace
    Eye tracking (2D)
      & Tobii Pro Glasses 3
      & 2D gaze, pupil diameter, validity
      & 50~Hz
      & 4 (one per P) \\
    \addlinespace
    Audio (close-talk)
      & DPA 4060
      & Mono, 48~kHz / 24-bit
      & 48~kHz
      & 4 (one per P) \\
    \addlinespace
    Audio (room)
      & DPA 4060
      & Mono, 48~kHz / 24-bit
      & 48~kHz
      & 1 \\
    \addlinespace
    Behavioural self-report
      & Android tablets
      & VAD probes, task responses
      & event-driven
      & 4 (one per P) \\
    \addlinespace
    Personality (pre-session)
      & Online questionnaire
      & BFI-44 item responses, demographics
      & one-shot
      & 4 (one per P) \\
    \addlinespace
    Conversation transcripts
      & WhisperX (per-channel)
      & Word-level diarised transcripts
      & event-driven
      & 1 (merged) \\
    \bottomrule
  \end{tabular}
\end{table}


\paragraph{Physiology.}
Each participant wears an EmotiBit sensor on the wrist, streaming
photoplethysmography, electrodermal activity, skin temperature, and inertial
motion.
Sampling rates are preserved at device-native resolution where possible.
% TODO: confirm channel-specific sampling rates.
Devices are mapped to seat identifiers during acquisition, and the mapping is
stored in the session metadata.

\paragraph{Eye tracking.}
Each participant wears a Tobii Pro Glasses 3 unit.
The present release includes 2D gaze, pupil diameter, validity flags, and
scene-camera timing information.
It does not include world-frame gaze alignment.
Therefore, direct claims about who looked at whom require additional processing
outside the current release.

\paragraph{Audio.}
Each participant has a close-talk lapel microphone recorded through a
multi-channel audio interface.
A room/ambient microphone may be used for noise and crosstalk
characterisation where available.
Close-talk audio is the basis for transcript generation and audio/prosodic
features.
Raw audio access is treated separately from lower-risk tabular modalities
because voice can be personally identifying.

\paragraph{Tablet self-report and task responses.}
Each participant uses a tablet for self-report probes and task-specific
responses.
Probe events, response timestamps, and task events are logged and aligned to
the common session timeline.

\paragraph{Synchronisation.}
Acquisition is organised around a common LSL timebase.
The processing pipeline uses LSL markers, session progress logs, and
event files to derive task windows and align modalities.
For the current release, the most important downstream artefacts are the
per-task windows, aligned events, and per-modality timestamped files.

%% ═══════════════════════════════════════════════════════════════════════════
\section{Processing and Release}
\label{sec:processing}

The release follows a BIDS-inspired structure \cite{gorgolewski2016brain}.
Each session contains per-modality subdirectories for physiological data,
eye tracking, audio-derived artefacts, behavioural/task responses, and
annotations.
The exact schema is documented in the release README and datasheet.

\paragraph{From raw acquisition to release files.}
The processing pipeline merges raw acquisition outputs, derives task windows
from LSL markers, writes canonical TSV/JSON sidecars, and produces
per-session quality-control summaries.
The intended design goal is deterministic conversion: re-running the pipeline
on the same raw inputs should produce the same release artefacts, aside from
non-substantive metadata timestamps.

\paragraph{Conversation transcripts.}
We aim to produce per-session transcripts from the close-talk audio channels.
Because each microphone is associated with a known seat identifier, channel
identity provides a strong prior for speaker assignment.
However, cross-talk between nearby lapel microphones remains a practical risk.
The transcript pipeline is therefore treated as a release artefact with
documented quality checks, not as unquestioned ground truth.

The planned transcript pipeline is:
\begin{enumerate}
  \item per-channel voice-activity detection,
  \item per-channel speech recognition,
  \item cross-talk suppression or quality flagging,
  \item merge into a single ordered transcript with P1--P4 speaker labels,
  \item name-redaction and basic privacy checks.
\end{enumerate}

If manual checking is completed before the internal freeze, we report
transcript-quality metrics in the appendix.
If not, the transcripts are released or described with explicit caveats rather
than used as validated analytical labels.

\paragraph{Anonymisation.}
Released files use anonymous seat identifiers rather than names.
Free-text fields and transcripts are scanned for direct identifiers and
redacted where necessary.
The real identity-to-seat mapping is excluded from the release.

\paragraph{Release tiers.}
The lower-risk tabular modalities are intended for public or open academic
release, subject to final consent and institutional checks.
Raw audio is handled as a separate access decision because of
voice re-identification risk.
The default conservative plan is to gate raw audio access under a Data Use
Agreement if the consent interpretation is ambiguous.
The paper will describe only the access structure that is actually in place at
submission time.

\paragraph{Hosting and metadata.}
The dataset is planned for release through a recognised ML/data hosting
platform, with Croissant metadata and Responsible AI fields.
A datasheet and release README document intended use, limitations, access
conditions, and maintenance policy.
% TODO: insert final anonymous review URL and version tag.

%% ═══════════════════════════════════════════════════════════════════════════
\section{Worked Example and Supported Analyses}
\label{sec:worked}

This section has two roles.
First, it provides one conservative worked example if the required feature
tables are ready and reliable.
Second, it specifies additional analyses that the dataset supports without
presenting them as complete benchmarks in this paper.

This is an intentionally conservative choice.
The dataset is multimodal and analytically rich, but the number of groups is
too small for leaderboard-style claims.

\subsection{Worked example: task-level multimodal characterisation}
\label{sec:worked_example}

% TODO: finalise after the Thursday checkpoint.
% The final question should be written before running the final analysis.

\paragraph{Question.}
We ask whether simple task-level summaries of physiology and/or audio-prosody
vary in interpretable ways across structured task phases and task outcomes.
The analysis is descriptive and intended to illustrate use of the release
structure.

\paragraph{Inputs.}
The preferred feature set combines:
(i) per-participant physiological summaries aggregated within each task window;
(ii) per-participant audio/prosodic summaries over the same windows; and
(iii) task-outcome descriptors from tablet responses.
If one modality is unavailable or unreliable, the worked example is run on the
available subset and the missing modality is reported as a limitation.

\paragraph{Outputs.}
The worked example reports a compact table or figure summarising the relation
between task phase, task outcome, and available multimodal features.
Effect sizes, confidence intervals, and descriptive patterns are preferred over
strong significance claims.

\paragraph{Interpretation.}
The worked example is a demonstration of dataset use, not a final claim about
group affect, personality, or performance.
Null or weak patterns are reported as such.

\subsection{Supported analysis: self-report dynamics}

The self-report probes support analysis of valence/arousal changes across
baseline and task phases.
Dominance ratings may be released where available, but dominance analysis is
not a core analysis in this submission.

\subsection{Supported analysis: personality and multimodal response}

The BFI-44 scores support exploratory analyses of whether stable personality
traits modulate physiological, self-report, or conversational responses across
tasks.
Given the sample size, such analyses should be treated as hypothesis-generating
rather than confirmatory.

\subsection{Supported analysis: task outcomes and interaction dynamics}

The structured tasks support outcome-oriented analyses, including
hidden-profile decision correctness, negotiation settlement, idea-generation
outputs, and public-goods contribution behaviour.
These outcomes can be linked to multimodal signals in follow-up work,
especially once transcript and audio-derived features are fully validated.

%% ═══════════════════════════════════════════════════════════════════════════
\section{Dataset Statistics and Quality}
\label{sec:stats}

\Cref{tab:stats} reports per-modality coverage across analysable sessions.
We report coverage per modality rather than requiring every session to be
complete for every possible analysis.

\begin{description}
  \item[Tier 1:]
  sessions with all core release modalities present for all four participants
  and no major documented anomalies.
  \item[Tier 2:]
  sessions with minor documented gaps, still usable for some analyses.
  \item[Tier 3:]
  sessions with significant gaps, included for transparency but excluded from
  the headline worked example unless specifically justified.
\end{description}

The worked example uses Tier~1 and eligible Tier~2 sessions for which the
required modalities are present.

%% Table: dataset_stats.tex
%% %% Table: dataset_stats.tex
%% %% Table: dataset_stats.tex
%% \input{tables/dataset_stats} from main.tex
%%
%% This table follows the audit in docs/data_audit.md. It does NOT use the
%% inflated "13 sessions / 19 hours" headline that appears in the ICMI draft.
%% Numbers below the per-modality coverage line are derived from QC artifacts
%% and remain TBD until the QC pipelines are re-run on the final session set.

\begin{table}[t]
  \centering
  \caption{AffectAI dataset statistics, reported by completeness tier.
    Tier 1 sessions are fully complete across all configured modalities for
    all four participants; Tier 2 has one minor documented gap;
    Tier 3 has multiple documented gaps and is excluded from the headline
    worked example in \cref{sec:worked}.
    Modality coverage is the percentage of sessions in the indicated tier
    set with usable data for that modality.
    QC metrics marked TBD will be filled once pipeline re-runs are
    complete; the underlying tools are part of this release.}
  \label{tab:stats}
  \small
  \begin{tabular}{lr}
    \toprule
    \textbf{Metric} & \textbf{Value} \\
    \midrule
    Total final-phase sessions                          & TBD \\
    \quad of which Tier 1 (fully complete)              & TBD \\
    \quad Tier 2 (nearly complete)                      & TBD \\
    \quad Tier 3 (significant documented gaps)          & TBD \\
    Unique groups                                       & TBD \\
    Unique participants                                 & TBD \\
    Total recorded duration                             & $\sim$TBD~h \\
    \midrule
    \multicolumn{2}{l}{\textit{Per-modality coverage (\% of Tier~1 $\cup$ Tier~2 sessions)}} \\
    \midrule
    EmotiBit (all four participants)                    & TBD\,\% \\
    Tobii 2D gaze (all four participants)               & TBD\,\% \\
    DPA close-talk audio (all four participants)        & TBD\,\% \\
    Room audio                                          & TBD\,\% \\
    VAD probes (all four participants, $\geq 90\%$ response rate)
                                                        & TBD\,\% \\
    BFI-44 personality (all four participants matched)  & TBD\,\% \\
    Conversation transcript                             & TBD\,\% \\
    \midrule
    \multicolumn{2}{l}{\textit{Synchronisation quality (Tier~1 $\cup$ Tier~2)}} \\
    \midrule
    Median cross-device LSL offset                      & TBD\,ms \\
    99\textsuperscript{th}-percentile LSL offset        & TBD\,ms \\
    Median frame-drop rate (close-talk audio)           & TBD\,\% \\
    \midrule
    \multicolumn{2}{l}{\textit{Transcript quality (if manually checked)}} \\
    \midrule
    Transcript QC status                                & TBD \\
    WER / diarisation note, if available                & TBD \\
    \bottomrule
  \end{tabular}
\end{table}
 from main.tex
%%
%% This table follows the audit in docs/data_audit.md. It does NOT use the
%% inflated "13 sessions / 19 hours" headline that appears in the ICMI draft.
%% Numbers below the per-modality coverage line are derived from QC artifacts
%% and remain TBD until the QC pipelines are re-run on the final session set.

\begin{table}[t]
  \centering
  \caption{AffectAI dataset statistics, reported by completeness tier.
    Tier 1 sessions are fully complete across all configured modalities for
    all four participants; Tier 2 has one minor documented gap;
    Tier 3 has multiple documented gaps and is excluded from the headline
    worked example in \cref{sec:worked}.
    Modality coverage is the percentage of sessions in the indicated tier
    set with usable data for that modality.
    QC metrics marked TBD will be filled once pipeline re-runs are
    complete; the underlying tools are part of this release.}
  \label{tab:stats}
  \small
  \begin{tabular}{lr}
    \toprule
    \textbf{Metric} & \textbf{Value} \\
    \midrule
    Total final-phase sessions                          & TBD \\
    \quad of which Tier 1 (fully complete)              & TBD \\
    \quad Tier 2 (nearly complete)                      & TBD \\
    \quad Tier 3 (significant documented gaps)          & TBD \\
    Unique groups                                       & TBD \\
    Unique participants                                 & TBD \\
    Total recorded duration                             & $\sim$TBD~h \\
    \midrule
    \multicolumn{2}{l}{\textit{Per-modality coverage (\% of Tier~1 $\cup$ Tier~2 sessions)}} \\
    \midrule
    EmotiBit (all four participants)                    & TBD\,\% \\
    Tobii 2D gaze (all four participants)               & TBD\,\% \\
    DPA close-talk audio (all four participants)        & TBD\,\% \\
    Room audio                                          & TBD\,\% \\
    VAD probes (all four participants, $\geq 90\%$ response rate)
                                                        & TBD\,\% \\
    BFI-44 personality (all four participants matched)  & TBD\,\% \\
    Conversation transcript                             & TBD\,\% \\
    \midrule
    \multicolumn{2}{l}{\textit{Synchronisation quality (Tier~1 $\cup$ Tier~2)}} \\
    \midrule
    Median cross-device LSL offset                      & TBD\,ms \\
    99\textsuperscript{th}-percentile LSL offset        & TBD\,ms \\
    Median frame-drop rate (close-talk audio)           & TBD\,\% \\
    \midrule
    \multicolumn{2}{l}{\textit{Transcript quality (if manually checked)}} \\
    \midrule
    Transcript QC status                                & TBD \\
    WER / diarisation note, if available                & TBD \\
    \bottomrule
  \end{tabular}
\end{table}
 from main.tex
%%
%% This table follows the audit in docs/data_audit.md. It does NOT use the
%% inflated "13 sessions / 19 hours" headline that appears in the ICMI draft.
%% Numbers below the per-modality coverage line are derived from QC artifacts
%% and remain TBD until the QC pipelines are re-run on the final session set.

\begin{table}[t]
  \centering
  \caption{AffectAI dataset statistics, reported by completeness tier.
    Tier 1 sessions are fully complete across all configured modalities for
    all four participants; Tier 2 has one minor documented gap;
    Tier 3 has multiple documented gaps and is excluded from the headline
    worked example in \cref{sec:worked}.
    Modality coverage is the percentage of sessions in the indicated tier
    set with usable data for that modality.
    QC metrics marked TBD will be filled once pipeline re-runs are
    complete; the underlying tools are part of this release.}
  \label{tab:stats}
  \small
  \begin{tabular}{lr}
    \toprule
    \textbf{Metric} & \textbf{Value} \\
    \midrule
    Total final-phase sessions                          & TBD \\
    \quad of which Tier 1 (fully complete)              & TBD \\
    \quad Tier 2 (nearly complete)                      & TBD \\
    \quad Tier 3 (significant documented gaps)          & TBD \\
    Unique groups                                       & TBD \\
    Unique participants                                 & TBD \\
    Total recorded duration                             & $\sim$TBD~h \\
    \midrule
    \multicolumn{2}{l}{\textit{Per-modality coverage (\% of Tier~1 $\cup$ Tier~2 sessions)}} \\
    \midrule
    EmotiBit (all four participants)                    & TBD\,\% \\
    Tobii 2D gaze (all four participants)               & TBD\,\% \\
    DPA close-talk audio (all four participants)        & TBD\,\% \\
    Room audio                                          & TBD\,\% \\
    VAD probes (all four participants, $\geq 90\%$ response rate)
                                                        & TBD\,\% \\
    BFI-44 personality (all four participants matched)  & TBD\,\% \\
    Conversation transcript                             & TBD\,\% \\
    \midrule
    \multicolumn{2}{l}{\textit{Synchronisation quality (Tier~1 $\cup$ Tier~2)}} \\
    \midrule
    Median cross-device LSL offset                      & TBD\,ms \\
    99\textsuperscript{th}-percentile LSL offset        & TBD\,ms \\
    Median frame-drop rate (close-talk audio)           & TBD\,\% \\
    \midrule
    \multicolumn{2}{l}{\textit{Transcript quality (if manually checked)}} \\
    \midrule
    Transcript QC status                                & TBD \\
    WER / diarisation note, if available                & TBD \\
    \bottomrule
  \end{tabular}
\end{table}


\paragraph{Synchronisation quality.}
% TODO: numbers from QC pipeline.
We report cross-device/session-level timing quality using LSL-derived
alignment summaries and per-session QC reports.
Where exact offsets cannot be estimated for a stream, this is reported as a
coverage limitation rather than imputed.

\paragraph{Modality coverage.}
Coverage is reported at session, participant, and modality level.
This avoids treating a partially missing modality as failure of the entire
session.

\paragraph{Audio and transcript quality.}
Audio quality is summarised using per-channel diagnostics where available.
Transcript quality is reported if manual checks are complete.
Otherwise, transcripts are described as automatically generated release
artefacts with quality caveats.

\paragraph{Self-report response rate.}
Probe response rates are reported overall and by session/task where available.
Sparse responses are treated as missingness, not interpolated into dense labels.

%% ═══════════════════════════════════════════════════════════════════════════
\section{Limitations}
\label{sec:limitations}

\paragraph{Sample size.}
The dataset contains \textbf{TODO} participants across \textbf{TODO} groups.
This is sufficient for dataset characterisation and exploratory multimodal
analysis, but not for broad benchmark or leaderboard claims.

\paragraph{Single site and language.}
All sessions were recorded at a single site in \textbf{[ANON]} and conducted
in English.
Turn-taking norms, humour, backchannels, and negotiation styles may therefore
reflect the local context.

\paragraph{Per-modality missingness.}
Not all sessions are complete for all modalities.
We report missingness explicitly and perform analyses only on sessions where
the required inputs are available.

\paragraph{Voice re-identification.}
Close-talk audio carries voice re-identification risk.
This motivates separate treatment of raw audio access and explicit restrictions
against speaker identification, voice cloning, and redistribution.

\paragraph{2D rather than world-frame gaze.}
Released gaze is 2D in the participant's own scene camera, not aligned to a
shared room coordinate system.
Cross-participant gaze-target analysis is therefore outside the current release.

\paragraph{Dominance and video-based analyses deferred.}
Although dominance-style self-report may exist in the raw/release data,
dominance analysis is deliberately deferred.
Similarly, multi-camera video, 3D pose, and world-frame gaze alignment are
reserved for follow-up work.

\paragraph{No clinical ground truth.}
The dataset is not a clinical affect dataset.
It should not be used for diagnosis, mental-health inference, personnel
evaluation, or surveillance-style deployment.

%% ═══════════════════════════════════════════════════════════════════════════
\section{Ethics, Access, and Responsible Use}
\label{sec:ethics}

\paragraph{Ethics approval and consent.}
The study was approved under \textbf{[ANON --- ethics reference]}.
Participants provided written informed consent before taking part.
The final release licence and access tiers follow the consent language and
institutional requirements.

\paragraph{Anonymisation.}
Released files use anonymous seat identifiers.
Direct identifiers are removed from tabular files and redacted from transcripts
where detected.
The identity mapping is not released.

\paragraph{Access policy.}
The dataset is intended for non-commercial academic research on multimodal
group interaction.
If raw audio is gated, access requires agreement to a Data Use Agreement.
The DUA prohibits speaker identification, voice cloning, workplace
surveillance, redistribution, and attempts to recover participant identity.

\paragraph{Out-of-scope use.}
The dataset is not appropriate for clinical diagnosis, individual scoring,
automated personnel evaluation, surveillance, or deployment in settings where
participants are monitored without informed consent.

\paragraph{Maintenance.}
The dataset will be maintained at \textbf{TODO hosting URL} for at least
\textbf{TODO} years from publication.
Errata, schema updates, and quality-control artefacts will be released under
versioned tags.
The version described in this paper is \textbf{TODO v1.0.0}.

%% ═══════════════════════════════════════════════════════════════════════════
%% Acknowledgements omitted for double-blind review.
%% \section*{Acknowledgements}

\bibliographystyle{plain}
\bibliography{references}

%% ── Appendix ──────────────────────────────────────────────────────────────
\appendix

\section{BFI-44 scoring and item list}
\label{app:bfi44}
% TODO: item list, reverse coding, scoring formula.

\section{Synchronisation pipeline detail}
\label{app:sync}
% TODO: LSL streams, task-window derivation, QC logic.

\section{Per-session quality table}
\label{app:quality}
% TODO: per-session × modality table.

\section{Transcription pipeline parameters}
\label{app:transcript}
% TODO: model versions, VAD parameters, cross-talk suppression, redaction.

\section{Datasheet for Datasets}
\label{app:datasheet}
% TODO: standard Datasheet sections.

\end{document}
